% Proved after four complete actual-file audits; assembled by root.
\subsection{A fixed detecting curve and its completed local sources}

\begin{lemma}[Proved: finite punctures give fixed local types on one curve]
\label{D:finite-puncture-detecting-curve-local-types}
Let $S$ be a normal proper integral surface over an algebraically
closed field $k$ of characteristic $p>0$, let $T$ be a finite closed
subset, and put $U=S\setminus T$ and $\Gamma=\pi_1(U)$.
There are a finite set $F\supset T$ of closed points of $S$, an
algebraically closed extension $\Omega/k$ of finite transcendence
degree, a smooth projective connected curve $C/\Omega$, a finite
nonempty marked set $P\subset C$, and a morphism
$C^\circ=C\setminus P\to U$ with the following properties.
\begin{enumerate}
\item The induced homomorphism $\pi_1(C^\circ)\to\Gamma$ is
surjective. Its target is the original group over $k$; no
surjectivity onto $\pi_1(U_\Omega)$ is asserted.
\item Put $R_s=\widehat{\mathcal O}_{S,s}$ and
$\Lambda_s=\pi_1(\operatorname{Spec}R_s\setminus\{\mathfrak m_s\})$
for $s\in F$. These are fixed complete excellent normal local
surface rings over the original field $k$. For every $x\in P$
there is an $s=s(x)\in F$ and a fixed morphism
\[
 \operatorname{Spec}\Omega((t_x))\longrightarrow
 \operatorname{Spec}R_s\setminus\{\mathfrak m_s\}
 \longrightarrow U .                                      \tag{D.DC.1}
\]
The corresponding homomorphisms of fundamental groups, with
compatible paths, give the local restriction of
$\pi_1(C^\circ)\to\Gamma$. Thus each full local representation
type at $x$ is a restriction of one representation of the fixed
group $\Lambda_s$, for every finite target group. If $s\notin T$,
the composite $\Lambda_s\to\Gamma$ is trivial.
\item There are fixed real numbers $b_x\ge0$, independent of the
finite quotient of $\Gamma$, which bound the highest classical
upper break of each induced local quotient at $x$. In particular,
if $L/\Omega((t_x))$ is its nontrivial Galois splitting extension,
with ramification index $e$ and different exponent $\delta$, then
\[
                    \delta/e\le b_x+1.                    \tag{D.DC.2}
\]
\end{enumerate}
All choices in this statement are independent of ranks, coefficient
fields and individual finite covers. Item~(2) concerns actual full
local types, not only their ramification breaks.
\end{lemma}

\begin{proof}
Choose the smooth projective surface $Y$ and SNC boundary $B$
used in Proposition~\ref{D:surface-curvewise-conductor} for the
pair $(S,T)$. Thus $Y\to S$ is proper birational, and after
enlarging a finite set $F\supset T$ its complement
$Y\setminus B$ maps isomorphically onto $S\setminus F$.
This uses normality in dimension two: a proper birational
morphism to $S$ is an isomorphism at the generic point of every
divisor, so its exceptional image is a finite set of closed points.
Chow's lemma and surface resolution produce this projective $Y$
even when $S$ itself is not projective. These are the same
choices and facts used in the cited proposition and in
Theorem~\ref{D:generic-surface-finite-punctures}.
If necessary add one smooth closed point to $F$ and blow it up;
this makes $B$ nonempty. The added exceptional curve corresponds
to a retained smooth point of $U$ and introduces no ramification
of a cover of $U$.

Fix a projective embedding $Y\hookrightarrow\mathbf P^N_k$.
Let $\Omega$ be an algebraic closure of
$k(z_0,\ldots,z_N)$, with the $z_i$ algebraically independent,
and let $C$ be the hyperplane section with these coefficients.
The geometric generic hyperplane is smooth and geometrically
irreducible. It avoids the finitely many nodes of $B$ and meets
every component of $B$ transversely. Indeed apply smooth Bertini
also to each smooth component of $B$; their zero-dimensional
sections are then reduced. Set $P=C\cap B_\Omega$. It is finite
and nonempty, since a hyperplane has positive intersection with
every projective curve. Smooth and irreducible Bertini here are
\citep[Lemmas~33.47.3 and~37.32.1--3, Tags~0FD6 and~0G4C,
with the written proofs]{StacksDGenericSurface}.

We check the simultaneous-cover quantifier explicitly. Let
$V\to U$ be any connected finite etale cover defined over $k$.
Its restriction $V_0$ to $S\setminus F$ is smooth and integral:
the total space $V$ is connected and normal, hence integral,
and a nonempty open of an integral scheme is integral.
Pull the hyperplane system back from $Y$ to $V_0$. It is
basepoint-free. At any closed point $v\in V_0$, etaleness and
the immersion in $\mathbf P^N$ allow two hyperplanes through
its image whose pullbacks have independent differentials at $v$.
Their intersection has a codimension-two component through $v$.
A third hyperplane avoiding the image of $v$ does not contain
that component. These are precisely conditions~(3) and~(4) of
Stacks, Lemma~37.32.3. Hence the geometric generic pullback
$V_0\times_{Y}C$ is geometrically irreducible, and is smooth
because it is etale over $C^\circ$. It is therefore connected.
The generic coefficients lie in the Bertini open for each
original cover separately. The one chosen generic point works
for all of them; no intersection of infinitely many open sets
at a closed point, or countability assumption, is used.
The connected-cover criterion proves item~(1).

For $x\in P$ let $s$ be its image in $F$. The map on local rings
from $\mathcal O_{S,s}$ to $\mathcal O_{C,x}$, followed by completion,
takes $\mathfrak m_s$ into $(t_x)$ in $\Omega[[t_x]]$.
It is therefore continuous and extends uniquely to
\[
                 R_s\longrightarrow\Omega[[t_x]].          \tag{D.DC.3}
\]
Its generic point does not map to the closed point of
$\operatorname{Spec}R_s$: the generic point of $C$ lies in
$Y\setminus B$, so some element of $\mathfrak m_s$ has nonzero
image in $\Omega[[t_x]]$, and hence invertible image in
$\Omega((t_x))$. This proves the first arrow in (D.DC.1).
The punctured spectrum of $R_s$ maps to $S\setminus T$.
Indeed the only point of $\operatorname{Spec}R_s$ above $s$
is its maximal ideal, and $\operatorname{Spec}\mathcal O_{S,s}$
has no other closed point of $S$ in its image. The composition
agrees with the punctured formal branch of $C^\circ\to U$,
by the uniqueness in (D.DC.3). This proves item~(2), including
compatibility for arbitrary finite torsors and their morphisms.
There is no appeal to invariance of punctured fundamental groups
under extending $k$, and no completed tensor product is needed.
The rings $R_s$ are normal domains by excellence and normality
of the finite-type local rings before completion. They have
coefficient and residue field $k$.
If $s\notin T$, a finite etale cover of $U$ restricts to one of
$\operatorname{Spec}R_s$. This strictly henselian local scheme
has only split finite etale covers, proving the last assertion.

For completeness we distinguish the coefficient characteristic
in the ramification argument. Apply the proof of
Proposition~\ref{D:surface-curvewise-conductor} over the fixed
algebraically closed field $\Omega$, to $(S_\Omega,T_\Omega)$
and the fixed SNC resolution $Y_\Omega$. It supplies a fixed
effective divisor $R$ on $Y_\Omega$ controlling all finite-free
systems of coefficient characteristic $\ell\ne p$ whose
monodromy comes from $\pi_1(U_\Omega)$. It is unnecessary
that its numerical coefficients equal those chosen over $k$.
Every original finite cover pulls back to such a cover over
$\Omega$. For its finite group $H$, take the faithful regular
representation $\mathbb F_\ell[H]$. The curvewise assertion
bounds its highest logarithmic slope at $x$ by
$b_x=\operatorname{mult}_x(C^*R)$, independently of $H$ and
its order. Over the perfect residue field $\Omega$, logarithmic
upper numbering is classical upper numbering. Faithfulness of
the regular representation means that every upper group above
$b_x$ in the induced quotient is trivial. This gives item~(3).
The curvewise comparison is exactly the SNC assertion of
\citet[Theorem~3.5, Remark~3.9 and Proposition~5.8(3)]
{C_HuTeyssierRam2025}; it uses the highest slope, not the
rank-weighted Swan conductor. Finally the elementary
different/Herbrand calculation in
Lemma~\ref{D:different-Frobenius-invariant-lattice} gives
(D.DC.2). In particular the bound uses the ramification index
of the curve's totally ramified local extension. No replacement
of a surface ramification index by a degree with a possibly
inseparable residue factor is made.
\end{proof}
